
\documentclass[12pt]{article}
\usepackage{amsmath}
\usepackage{amsfonts}
\usepackage{amssymb}
\usepackage{geometry}
\geometry{a4paper, margin=0.8in}
\pagenumbering{gobble}

\title{Homework solutions for 02YELMA \\ Chapter: Mathematical tools}
\begin{document}
\date{}
\maketitle
\vspace{-0.8in}

\begin{enumerate}

\item 
\[
f(x,y,z)=xyz
\]
\[
\nabla f = (yz,\, xz,\, xy)
\]
\[
\nabla f(1,1,1) = (1,1,1), \quad |\nabla f|=\sqrt{3}
\]

The gradient points in the direction of maximum increase of the function, and its magnitude gives the maximum rate of change. The directional derivative in any unit direction $\hat{u}$ is $\nabla f \cdot \hat{u}$.

\item 
\[
\vec{F}=(x^2,-y^2,z^2)
\]
\[
\nabla \cdot \vec{F} = 2x - 2y + 2z
\]
\[
\nabla \cdot \vec{F}(1,-1,1)=6
\]

A positive divergence indicates that the point behaves like a source, meaning the field locally flows outward from that point.

\item 
\[
\vec{F}=(y,z,x)
\]
\[
\nabla \times \vec{F} = (-1,-1,-1)
\]

The curl represents local rotation. A constant nonzero curl indicates uniform rotational behavior, with the direction of the curl vector giving the axis of rotation.

\item 
\[
\nabla \times (\nabla T)=0
\]
This follows from equality of mixed partial derivatives.

\[
\nabla \cdot (\nabla \times \vec{T})=0
\]
This also follows from cancellation of mixed partial derivatives. See Griffiths, Introduction to Electrodynamics, 3th edition, section 1.2.7 for more details.

\item For example, let $\vec{A}=\vec{B}$ and let $\vec{C}$ be perpendicular to both. Then the directions of $(\vec{A}\times \vec{B})\times \vec{C}$ and $\vec{A}\times (\vec{B}\times \vec{C})$ would not be the same, so the triple product is not associative in general.

\item 
\[
\vec{F}=(x,y)
\]
\[
\oint \vec{F}\cdot d\vec{r} = 0
\]

The field is conservative since $\vec{F}=\nabla\left(\frac{x^2+y^2}{2}\right)$, so the line integral over any closed path is zero. Or, if you really want to compute it:

\[
\vec{F}=(x,y), \qquad \oint \vec{F}\cdot d\vec{r}
= \int_{AB} + \int_{BC} + \int_{CA}
\]

\textbf{AB:} \; $x=t$, $y=1$, $t\in[0,1]$
\[
\int_{AB} \vec{F}\cdot d\vec{r}
= \int_0^1 t\,dt = \frac{1}{2}
\]

\textbf{BC:} \; $x=1-t$, $y=1-t$, $t\in[0,1]$
\[
\int_{BC} \vec{F}\cdot d\vec{r}
= \int_0^1 -2(1-t)\,dt = -1
\]

\textbf{CA:} \; $x=0$, $y=t$, $t\in[0,1]$
\[
\int_{CA} \vec{F}\cdot d\vec{r}
= \int_0^1 t\,dt = \frac{1}{2}
\]

\[
\oint \vec{F}\cdot d\vec{r}
= \frac{1}{2} - 1 + \frac{1}{2} = 0
\]

\[
\vec{F} = \nabla\!\left(\frac{x^2+y^2}{2}\right)
\;\Rightarrow\;
\oint \vec{F}\cdot d\vec{r} = 0
\]


\item 
\[
\vec{F}=(x,y,z)=\vec{r}
\]
\[
\iint \vec{F}\cdot d\vec{A} = 4\pi
\]

This gives the total outward flux through the sphere, representing the net strength of a source at the origin.

\item 
\[
\rho = x+y+z
\]
\[
M = \int_0^3 \int_0^{2\pi} \int_0^1 (r\cos\theta + r\sin\theta + z)\, r\,dr\,d\theta\,dz
\]
\[
M = \frac{9\pi}{2}
\]

The terms involving $\cos\theta$ and $\sin\theta$ vanish due to symmetry, leaving only the contribution from $z$.

\item 
\[
\vec{F}=(x^2,y^2,z^2)
\]
\[
\nabla \cdot \vec{F} = 2x+2y+2z
\]
\[
\iiint (2x+2y+2z)\,dV = 3
\]

The Divergence Theorem converts the surface flux into a volume integral of the divergence.

\item 
\[
\vec{F}=(-y,x,z)
\]
\[
\nabla \times \vec{F} = (0,0,2)
\]
\[
\oint \vec{F}\cdot d\vec{r} = 2\pi
\]

Stokes' Theorem relates circulation around the boundary to the flux of the curl through the surface.

\item 
\[
\vec{V} = k r^{-p}\hat{r}
\]
\[
\nabla \cdot \vec{V} = k r r^{-p} \left(2 - p\right) = 0
\]
This condition is satisfied for all $r\neq 0$ if and only if $p=2$. For $p=2$, the flux is independent of radius, meaning no sources or sinks exist away from the origin. The origin itself is a point source with strength $4\pi k$.

\item 
\[
R = \sqrt{(x-x')^2+(y-y')^2+(z-z')^2}
\]
\[
\frac{\partial}{\partial x'}\left(\frac{1}{R}\right)
= \frac{x-x'}{R^3}, \quad
\frac{\partial}{\partial y'}\left(\frac{1}{R}\right)
= \frac{y-y'}{R^3}, \quad
\frac{\partial}{\partial z'}\left(\frac{1}{R}\right)
= \frac{z-z'}{R^3}
\]
\[
\nabla' \left(\frac{1}{R}\right)
= \frac{\vec{r}-\vec{r}'}{R^3}
= \frac{\vec{R}}{R^3}
= \frac{\hat{R}}{R^2}
\]

The positive sign arises because differentiation is taken with respect to the primed coordinates.

\end{enumerate}


\end{document}
